\documentclass[conference]{IEEEtran}
\IEEEoverridecommandlockouts
% The preceding line is only needed to identify funding in the first footnote. If that is unneeded, please comment it out.
\usepackage{cite}
\usepackage{amsmath,amssymb,amsfonts}
\usepackage{algorithmic}
\usepackage{graphicx}
\usepackage{textcomp}
\usepackage{xcolor}
\def\BibTeX{{\rm B\kern-.05em{\sc i\kern-.025em b}\kern-.08em
    T\kern-.1667em\lower.7ex\hbox{E}\kern-.125emX}}
\begin{document}

\title{Motion-Adaptive Internal Force Allocation for Dual-Arm Cooperative Manipulation // Or // Task-on-Demand Grasp Force Optimization for Dual-Arm Object Transport\\

}

\author{\IEEEauthorblockN{1\textsuperscript{st} Luca Morello}
\IEEEauthorblockA{\textit{dept. name of organization (of Aff.)} \\
\textit{name of organization (of Aff.)}\\
City, Country \\
email address or ORCID}
\and
\IEEEauthorblockN{2\textsuperscript{nd} Given Name Surname}
\IEEEauthorblockA{\textit{dept. name of organization (of Aff.)} \\
\textit{name of organization (of Aff.)}\\
City, Country \\
email address or ORCID}
\and
\IEEEauthorblockN{3\textsuperscript{rd} Given Name Surname}
\IEEEauthorblockA{\textit{dept. name of organization (of Aff.)} \\
\textit{name of organization (of Aff.)}\\
City, Country \\
email address or ORCID}
}

\maketitle

\begin{abstract}
ToDo
\end{abstract}

\begin{IEEEkeywords}
ToDo \cite{uchiyama1988symmetric}
\end{IEEEkeywords}

\section{Introduction}
\begin{itemize}
    \item Dual-arm cooperative manipulation needs enough squeeze force to prevent slip, but excess squeeze force is wasted effort / risks damaging or destabilizing fragile objects.
    \item Human grip-force control scales force with load/motion, not statically (Westling \& Johansson 1984)
    \item Most cooperative-manipulation force-distribution work optimizes internal force *statically or per-instant*; most motion-adaptive grip-force work is single-gripper/slip-triggered, not posed as a dual-arm allocation problem.
    \item \textbf{Gap}: no existing framework couples 
    \begin{enumerate}
        \item nullspace-based internal/external decomposition,
        \item real-time QP allocation between two arms
        \item motion-derived force demand that ramps up specifically during object transport and relaxes at rest.
    \end{enumerate}
    \item Contributions
    \begin{enumerate}
        \item A QP formulation for real-time squeeze-force allocation between two cooperating arms with friction-aware bounds.
        \item A motion-adaptive demand term that couples the desired internal force level to the object's trajectory state.
        \item Hardware/sim validation on a dual xArm7 platform (mc\_rtc).
    \end{enumerate}
\end{itemize}



\section{Related Work}
\begin{itemize}
    \item A. Internal/external force decomposition and QP-based load distribution in dual-arm systems.
    \item B. Friction-cone-constrained grasp/contact force optimization (exact SOCP vs. linear/QP approximations).
    \item C. Motion- and slip-adaptive grip force control (bio-inspired and robotic).
    \item D. (brief) Impact-aware contact establishment — cited as context for your grasp-initiation phase, explicitly *not* claimed as a contribution here.
\end{itemize}


%\section{System Overview and Problem Formulation}
%- FSM overview (Idle → Independent → Collaborative: Build-Grasp → Trajectory → Cooperative Motion) — one figure.
%- Grasp matrix `G`, nullspace projector `Pint = I - G⁺G`, internal wrench definition.
%- Definition of the squeeze direction `n_squeeze` and its projection through `Pint`.
%- Problem statement: find per-arm normal forces `x = [F_L, F_R]` that track a desired internal force level while respecting friction and actuation bounds, at every control cycle.





%\section{Motion-Adaptive Force Demand ("Task on Demand")}
%- Definition of `F_demand(t_norm)` and its physical interpretation (feedforward margin that peaks during transport, vanishes at rest/at trajectory endpoints).
%- Design rationale vs. constant-safety-margin baselines; relation to human grip-force/load-force coupling.
%- **(If implemented per 4.2)**: extension to a velocity/acceleration-driven demand term, and comparison with the trajectory-phase-scheduled version.



%\section{QP-Based Squeeze Force Optimization}
%- Derivation of `H`, `g` from `n_squeeze`, `Pint`, `S_n`, `w_fixed` (as in `buildQPProblem`).
%- Role of `α` (minimization gain) vs. `β` (inter-arm balance gain) — what each term trades off, with a small sensitivity figure/table.
%- Friction-aware bound construction (`computeBounds`) and explicit discussion of the linearized, one-cycle-lagged approximation vs. an exact SOCP friction cone — cite Buss/Hashimoto/Moore and the 2026 SOCP-vs-QP paper here.
%- Real-time solve (QLD), computational cost note (2-variable QP — should be near-instant; report actual solve time).



%\section{Grasp Establishment (brief, supporting-only)}
%- One paragraph + one figure: reflected-mass-scaled virtual mass during `BUILD_GRASP` to soften first contact.
%- Explicitly framed as adopting an existing idea (reflected mass / preemptive impedance softening — cite Khatib, Arita), not as a contribution of this paper.
%- State current limitation honestly (e.g., "this is a heuristic softening term; a full impact-aware treatment as in [van Steen et al., Wang et al.] is left as future work").



%\section{Experimental Validation}
%- Platform: dual xArm7, `mc_rtc`.
%- Protocol: pick-carry-place trajectory, with/without motion-adaptive demand term (ablation), possibly varying object mass/friction.
%- Metrics: internal force tracking error vs. desired `λ_desired`; peak/RMS tangential force relative to friction bound (slip margin proxy); total squeeze effort (∫F dt) as an efficiency metric; success/failure under a perturbation or higher-speed trajectory.
%- Figures: (1) FSM/system diagram, (2) internal force vs. trajectory phase for both conditions, (3) friction-margin utilization, (4) solve-time histogram.


%\section{Discussion and Limitations}
%- Nonuniqueness of the internal/external decomposition (cite §3) — your α/β weighting is a design choice, not "the" optimum; discuss what it privileges.
%- Linearized, lagged friction bound vs. exact cone — when it could fail (fast tangential force changes).
%- Trajectory-phase-scheduled demand vs. true velocity-feedback demand — generalizability limits if you keep the scheduled version.
%- Grasp-establishment heuristic is not yet a full impact-aware controller.


%\section{Conclusion and Future Work}
%- Summary of contribution and results.
%- Future work: (1) closed-loop velocity/acceleration-driven demand, (2) exact/SOCP friction handling if solve time allows, (3) completing the impact-aware contact establishment as a follow-up paper, explicitly pointing to the reference-spreading / impact-aware QP literature as the target formalism.


\begin{table}[htbp]
\caption{Table Type Styles}
\begin{center}
\begin{tabular}{|c|c|c|c|}
\hline
\textbf{Table}&\multicolumn{3}{|c|}{\textbf{Table Column Head}} \\
\cline{2-4} 
\textbf{Head} & \textbf{\textit{Table column subhead}}& \textbf{\textit{Subhead}}& \textbf{\textit{Subhead}} \\
\hline
copy& More table copy$^{\mathrm{a}}$& &  \\
\hline
\multicolumn{4}{l}{$^{\mathrm{a}}$Sample of a Table footnote.}
\end{tabular}
\label{tab1}
\end{center}
\end{table}



\section*{Acknowledgment}


\section*{References}


\bibliographystyle{IEEEtran}
\bibliography{references}

\end{document}


